Los siguientes puntos describen las bases de datos con las cuales los análisis econométricos propuestos pueden ser implementados. Se trata de bases de datos que se encuentran disponibles dentro de YPF. 

\begin{itemize}
    \item \emph{DIMENSIONAL ESTACIONES DE SERVICIO / TIENDAS}. Datos de ubicación de cada una de las estaciones de servicio: provincia, localidad, coordenadas geográficas, nombre, detalle de red (propia, terceros, ACA), etc. Eventualmente, también se cuenta con una base con datos más demográficos de la zona (NSE, densidad poblacional).
    \item \emph{TRANSACCIONAL ESTACIONES DE SERVICIO}. Datos agrupados a nivel de despacho (transacción), con información sobre tipo de combustible, monto de la carga (litros y \$), fecha y hora, estación en la que se realizó, surtidor, etc.
    \item \emph{TRANSACCIONAL TIENDAS}. Datos agrupados a nivel de transacción, con detalle de producto, categoría, monto comprado, fecha y hora, estación, etc.
    \item \emph{TRANSACCIONAL SERVICLUB}. Datos agrupados a nivel de transacción, con detalle de cliente, combustible, monto (litros y \$), fecha y hora y medio de pago (APP, MP, MODO o desconocido). Para las compras en la tienda Full, solo se dispone de detalle de producto para aquellos clientes que compran con la app.
\end{itemize}





